% Source: lecture_notes_in deep_learning_ch1_summary.tex
\begin{notechapter}
\chapter{Deep Learning: An Overview}\label{ch:dl-lecture-notes-ch1}
\chaptersubtitle{A summary of Chapter 1 of \textit{Lecture Notes in Deep Learning}}

\section{Introduction to Deep Learning}
Deep learning is built on the idea of composing simple functions (layers) to form complex models, reminiscent of Kolmogorov's solution to Hilbert's 13th problem. Its recent success is attributed to:
\begin{itemize}
    \item Multi-layer architectures that extract hierarchical features (e.g., AlexNet)
    \item Availability of large datasets (``big data'')
    \item Advances in parallel computing (GPUs)
\end{itemize}

\section{From Linear Models to Neural Networks}
The goal of machine learning is to learn a model $f: \mathcal{X} \to \mathcal{Y}$ from data. A linear model is defined as:
\[
f(\mathbf{x}) = \mathbf{w}^\top \mathbf{x} + b
\]
However, linear models are limited—they cannot represent nonlinear relationships like the \textbf{XOR} function. To overcome this, we introduce \textbf{nonlinear activation functions} $\sigma$:
\[
f(\mathbf{x}) = \sigma(\mathbf{w}^\top \mathbf{x} + b)
\]
Common activation functions include:
\begin{itemize}
    \item \textbf{ReLU}: $\text{ReLU}(x) = \max(0, x)$
    \item \textbf{Sigmoid}: $\sigma(x) = \frac{1}{1 + e^{-x}}$
    \item \textbf{tanh}: $\tanh(x) = \frac{e^x - e^{-x}}{e^x + e^{-x}}$
    \item \textbf{Leaky ReLU}, \textbf{GELU}, etc.
\end{itemize}

\section{Neural Network Architecture}
By cascading multiple layers, we form a \textbf{neural network}:
\[
\mathbf{h}_l = \sigma(\mathbf{W}_l^\top \mathbf{h}_{l-1} + \mathbf{b}_l), \quad l = 1, \dots, L
\]
with $\mathbf{h}_0 = \mathbf{x}$ (input) and $\mathbf{h}_L = \mathbf{z}$ (output). Such a network is called a:
\begin{itemize}
    \item \textbf{Multi-Layer Perceptron (MLP)}
    \item \textbf{Fully Connected Network (FCN)}
    \item \textbf{Deep Neural Network (DNN)} if $L$ is large
\end{itemize}

\section{Training Neural Networks}
Training involves minimizing a \textbf{loss function} $\ell(\theta)$ over parameters $\theta = \{\mathbf{W}_l, \mathbf{b}_l\}$.

\subsection{Theoretical Foundation}
The ideal loss is the \textbf{expected risk} over the data distribution $\mathcal{D}$:
\[
\ell(\theta) = \mathbb{E}_{\mathbf{x} \sim \mathcal{D}} \left[ \mathcal{E}(\mathbf{x}; \theta) \right]
\]
where $\mathcal{E}(\mathbf{x}; \theta)$ is the error for a single data point. For MSE with ground truth $y(\mathbf{x})$:
\[
\ell(\theta) = \mathbb{E}_{\mathbf{x} \sim \mathcal{D}} \left[ \| y(\mathbf{x}) - f(\mathbf{x}, \theta) \|^2 \right]
\]

\subsection{Gradient Descent Optimization}
The optimal update direction is given by the \textbf{expected gradient}:
\[
\nabla_\theta \ell(\theta) = \mathbb{E}_{\mathbf{x} \sim \mathcal{D}} \left[ \nabla_\theta \mathcal{E}(\mathbf{x}; \theta) \right]
\]
Parameters are updated via gradient descent:
\[
\theta^{(t+1)} = \theta^{(t)} - \eta \nabla_\theta \ell(\theta^{(t)})
\]
where $\eta$ is the \textbf{learning rate}.

\subsection{Practical Gradient Approximations}
Since the true distribution $\mathcal{D}$ is unknown, we approximate the expectation:

\begin{itemize}
    \item \textbf{Batch Gradient Descent}:
    \[
    \mathbb{E}_{\mathbf{x} \sim \mathcal{D}}[\cdot] \approx \frac{1}{N} \sum_{i=1}^N [\cdot]
    \]
    Uses full dataset: stable but computationally expensive.
    
    \item \textbf{Stochastic Gradient Descent (SGD)}:
    \[
    \mathbb{E}_{\mathbf{x} \sim \mathcal{D}}[\cdot] \approx \nabla_\theta \mathcal{E}(\mathbf{x}_i; \theta)
    \]
    Uses one random sample: noisy but efficient (called ``online'' in the book).
    
    \item \textbf{Mini-batch SGD}:
    \[
    \mathbb{E}_{\mathbf{x} \sim \mathcal{D}}[\cdot] \approx \frac{1}{|\mathbb{B}|} \sum_{i \in \mathbb{B}} \nabla_\theta \mathcal{E}(\mathbf{x}_i; \theta)
    \]
    Uses a random subset: balances noise and efficiency (the book calls this ``SGD'').
\end{itemize}

\subsection{Empirical Loss Functions}
\begin{itemize}
    \item \textbf{Mean Squared Error (MSE)} for regression:
    \[
    \ell(\theta; \mathbb{X}) = \frac{1}{N} \sum_{i=1}^N \| \mathbf{y}_i - f(\mathbf{x}_i, \theta) \|^2
    \]
    \item \textbf{Cross-Entropy (CE)} for classification (with softmax):
    \[
    \mathcal{E}_{CE} = -\sum_{i=1}^{C} y_i \log(\tilde{z}_i), \quad \tilde{z}_i = \frac{\exp(z_i)}{\sum_j \exp(z_j)}
    \]
\end{itemize}

\subsection{Training Process}
\begin{itemize}
    \item One \textbf{epoch} = one pass through the entire training data
    \item Learning rate selection via \textbf{validation}
    \item Use of \textbf{schedulers} (e.g., 1-cycle) to adapt $\eta$ during training
\end{itemize}

\section{PyTorch Implementation}
The chapter introduces \textbf{PyTorch} for building and training neural networks. Key steps include:
\begin{itemize}
    \item Defining network architecture via \texttt{nn.Module}
    \item Using \texttt{DataLoader} for mini-batch processing
    \item Implementing training and test loops with backpropagation
    \item Example: Image classification on \textbf{Fashion-MNIST}
\end{itemize}

\section{Key Insights}
\begin{itemize}
    \item Depth enables hierarchical feature learning
    \item Nonlinear activations are essential for expressivity
    \item The core optimization minimizes expected risk via gradient approximations
    \item Different gradient approximations offer trade-offs between stability and efficiency
    \item Frameworks like PyTorch abstract away differentiation and hardware details
\end{itemize}
\end{notechapter}
